\subsection{Temporal reachability, counting, and cut problems}

Molter, Renken, and Zschoche~\citep[Theorem 14]{molter2024} give
polynomial-time algorithms for minimizing temporal reachability by deleting
or delaying time-edges when the underlying graph is a forest. Their
Lemma 2 transfers deletion to delay while preserving the forest property.
The stronger weighted forest delay formulation
\citep[Theorem 15, pp.~10--12]{molter2024} permits nonnegative weights on
reached vertices and undelayable time-edges, with a cardinality bound on
the modified time-edges. Its dynamic program records a vertex, a delay
budget, its first arrival time, and whether that arrival is supplied from
inside or outside its rooted subtree; a table entry minimizes the total
weight of reached vertices. The underlying temporal graph is undirected,
and the forest result allows arbitrary nonnegative traversal times,
including the strict convention with traversal time one. Thus tree
structure, strictness, and protected contacts alone are not novelty
distinctions, and the source's vertex weights must not be confused with
arbitrary event-deletion costs.

The final journal version also gives a $4^x|G|^{O(1)}$ algorithm for both
reachability problems in the non-strict setting, where $x$ is the timed
feedback vertex number~\citep[Theorem 16]{molter2024}. This parameter is
zero on underlying forests. The output contract here instead requires
minimum original-event deletion cost at every attainable exact pair of
root-ending two- and three-event tuple counts, subject to the preservation
quota and shared event decisions. A derivation from the reachability
results must therefore supply an encoding or sufficient state for these
multiplicities, costs, and constraints, as well as the claimed dependence
on preservation loss. Appendix~\ref{app:reachability} gives a direct-summary
counterexample: the same reached sets and earliest arrivals can coexist
with different tuple counts. It does not exclude richer states or auxiliary
reductions. The journal article extends the earlier conference version
\citep{molter2021}. The comparison above uses the final journal
statements and proofs, with their journal theorem numbering.

Enright, Meeks, and Molter~\citep[Theorem 2]{enright2025counting} count temporal
paths on forests in $O(|V|\tau^2)$ time, where $\tau$ is the lifetime. Their
Theorems 6--7 treat further structural parameters. Fixed-deletion counting
does transfer to the present setting: Appendix~\ref{app:counting} gives a single
normalization preserving strict ties, event identities, and both counts
under every deletion subset, with lifetime at most $M$. Thus strictness,
orientation, and parallel contacts do not obstruct that counting transfer.
The additional task here is to choose one shared deletion set and compute
minimum costs for the entire attainable profile with the stated dependence
on preservation loss.

Robust Temporal Cut~\citep{enright2026robust} removes contacts so that no
source--target temporal path survives even after bounded timestamp
perturbations. For undirected temporal graphs, its stated tractability
parameters include timed vertex
cover, vertex cover plus solution size, and temporal neighborhood diversity
plus solution size. This is a related temporal-deletion problem with a
different feasibility condition. The existence of those algorithms does not
itself supply an exact reduction for two simultaneous multiplicity counts.
Any proposed transfer must also preserve its controlling parameters.

\subsection{Cover formulations and the parameters they retain}

There is a direct exact set-system representation. Each baseline two-walk is
a red element, each baseline three-walk is a blue element, and each optional
event supplies the set of elements containing its original identifier.
Selecting event sets incurs their deletion costs. A profile entry
$C(\ell,w)$ requires red union coverage exactly $\ell$ and blue union
coverage exactly $N_3-w$. An upper quota on surviving triples instead
requires a partial blue target. Shared incidences matter: a walk destroyed
by several deletions is still counted once.

This representation relates the problem to red--blue set cover. Abidha and
Ashok~\citep[Theorem 6 in the author preprint]{abidha2024} give an exact
fixed-dimensional hyperplane algorithm parameterized by red coverage for
covering all blue points. The temporal witness incidence admits a
three-dimensional coordinate-plane representation, but the compared
feasible family has no independent charge for selected planes and freely
selects red-free sets. First-event choices in the temporal model can be red-free and still
cost money. The partial-blue approximation setting of
Chlamt\'a\v{c}, Makarychev, and Vakilian~\citep[Appendix A]{chlamtac2023}
is another close precedent for the benefit--damage formulation; its
approximation guarantee does not by itself decide exact profile feasibility
and minimum event cost.

Madireddy, Nandy, and Pandit~\citep[Theorems 1--2]{madireddy2025} give exact
algorithms for intervals and for a restricted axis-parallel-line case. The
direct temporal witness representation need not be an interval system.
With one first event $a$, two middle events $b_1,b_2$, and two final events
$c_1,c_2$, all temporally compatible, the four triples $q_{ij}=(a,b_i,c_j)$
induce the event sets
\[
 \{q_{11},q_{12}\},\quad \{q_{21},q_{22}\},\quad
 \{q_{11},q_{21}\},\quad \{q_{12},q_{22}\}.
\]
An interval order would need four distinct adjacent pairs on four blue
points, which have only three adjacencies. The restricted line algorithm
also relies on disjoint red incidences, whereas a temporal red pair belongs
to both its middle and final events. These observations concern the direct
representation and do not rule out another exact encoding.

The bipartite partial-cover framework of
Mkrtchyan and coauthors~\citep{mkrtchyan2024} includes a selected-vertex cost
budget and covered-edge profit targets. Their hardness theorem with two side budgets
provides the starting problem for the reduction in Appendix~A. For the positive comparisons,
neither selected cost nor covered profit is inherently bounded by the preservation loss $R$.
For example, take $n$ optional unit-cost first events before one retained
middle/final pair and require preservation of that pair. Then $R=0$ while
destroying all triples requires $n$ deletions. A target leaving at most
$\lfloor n/2\rfloor$ triples makes both the destroyed-triple target and the
surviving-triple allowance grow at $R=0$. An algorithm controlled by those
quantities therefore needs an additional parameter argument to yield the bound in Theorem~\ref{thm:main}. This is a statement about the indicated transfer, not about all
possible reductions.
